
\documentclass{article}
\usepackage{colortbl}
\usepackage{makecell}
\usepackage{multirow}
\usepackage{supertabular}

\begin{document}

\newcounter{utterance}

\centering \large Interaction Transcript for game `wordle', experiment `medium\_frequency\_words\_no\_clue\_no\_critic', episode 10 with gemma{-}4{-}26b{-}a4b{-}it.
\vspace{24pt}

{ \footnotesize  \setcounter{utterance}{1}
\setlength{\tabcolsep}{0pt}
\begin{supertabular}{c@{$\;$}|p{.15\linewidth}@{}p{.15\linewidth}p{.15\linewidth}p{.15\linewidth}p{.15\linewidth}p{.15\linewidth}}
    \# & \multicolumn{2}{c}{Player} && \multicolumn{2}{c}{Game Master} \\
    \hline

    \theutterance \stepcounter{utterance}  
    & & & \multicolumn{4}{p{0.6\linewidth}}{
        \cellcolor[rgb]{0.9,0.9,0.9}{
            \makecell[{{p{\linewidth}}}]{
                \texttt{\tiny{[P1$\langle$GM]}}
                \texttt{Tu esi kalbos burtininkas, kuriam patinka spėti žodžius, laikantis pateiktų taisyklių.} \\
\\ 
\texttt{Sveiki atvykę į Wordle! Turite šešis bandymus atspėti tikslinį žodį – galiojantį penkių mažųjų raidžių anglų kalbos žodį. Prašome naudoti žymas „paaiškinimas:" ir „spėjimas:", kad pateiktumėte glaustą paaiškinimą kiekvienam spėjimui.} \\
\\ 
\texttt{Pavyzdžiui, jei jūsų spėjimas yra „įmonė", jūsų atsakymas turėtų būti} \\
\texttt{paaiškinimas: tai yra dažnas penkių raidžių anglų kalbos žodis, ir savo spėjimą pradedu šiuo žodžiu.} \\
\texttt{spėjimas: įmonė} \\
\\ 
\texttt{Po kiekvieno spėjimo jūsų atsakymas bus patikrintas ir gausite grįžtamąjį ryšį, nurodantį, kurios raidės yra teisingos (žalios), kurios raidės teisingos, bet netinkamoje pozicijoje (geltonos), ir kurios raidės neteisingos (raudonos). Šis grįžtamasis ryšys gali būti naudingas sprendžiant, kurias raides įtraukti arba atmesti kitame spėjime.} \\
\\ 
\texttt{Pavyzdžiui, grįžtamasis ryšys žodžiui „įmonė" galėtų būti:} \\
\texttt{spėjimo\_atsiliepimas: į<geltona> m<geltona> o<žalia> n<geltona> ė<raudona>} \\
\\ 
\texttt{Paaiškinime turi būti detalių apie tai, kaip spėjimo\_atsiliepimas naudojamas naujam spėjimui suformuoti.} \\
\\ 
\texttt{Negeneruokite jokio kito teksto.} \\
\\ 
\texttt{Pradėkime nuo jūsų pirmojo spėjimo.} \\
            }
        }
    }
     \\ \\

    \theutterance \stepcounter{utterance}  
    & \multicolumn{4}{p{0.6\linewidth}}{
        \cellcolor[rgb]{0.561,0.737,0.561}{
            \makecell[{{p{\linewidth}}}]{
                \texttt{\tiny{[P1$\rangle$GM]}}
                \texttt{paaiškinimas: Pradžiai renkuosi dažną anglų kalbos žodį su dažnomis raidėmis (A, E, I, O, U ir consonantes), kad galėčiau greitai nustatyti pagrindines balses bei raidžių buvimą.} \\
\texttt{spėjimas: audio} \\
            }
        }
    }
    & & \\ \\

    \theutterance \stepcounter{utterance}  
    & & & \multicolumn{2}{p{0.3\linewidth}}{
        \cellcolor[rgb]{0.9,0.9,0.9}{
            \makecell[{{p{\linewidth}}}]{
                \texttt{\tiny{[GM$|$GM]}}
                \texttt{attempts = 1} \\
\texttt{target\_word = akmuo} \\
\texttt{guess = audio} \\
\texttt{guess\_feedback = a<green> u<yellow> d<red> i<red> o<green>} \\
            }
        }
    }
    & & \\ \\

    \theutterance \stepcounter{utterance}  
    & & & \multicolumn{4}{p{0.6\linewidth}}{
        \cellcolor[rgb]{0.9,0.9,0.9}{
            \makecell[{{p{\linewidth}}}]{
                \texttt{\tiny{[P1$\langle$GM]}}
                \texttt{spėjimo\_atsiliepimas: a<žalia> u<geltona> d<raudona> i<raudona> o<žalia>} \\
\\ 
\texttt{Atsakymą pateikite tik šiuo formatu.} \\
\texttt{paaiškinimas: detalės} \\
\texttt{spėjimas: žodis} \\
            }
        }
    }
     \\ \\

    \theutterance \stepcounter{utterance}  
    & \multicolumn{4}{p{0.6\linewidth}}{
        \cellcolor[rgb]{0.561,0.737,0.561}{
            \makecell[{{p{\linewidth}}}]{
                \texttt{\tiny{[P1$\rangle$GM]}}
                \texttt{paaiškinimas: Žinau, kad žodis prasideda su „a“ ir baigiasi su „o“. Raidė „u“ yra žodžio dalyje, bet ne trečioje pozicijoje. Raidės „d“ ir „i“ yra neteisingos. Dabar ieškau žodžio, kuriame būtų „u“ (ne trečioje vietoje) ir naujų raidžių, pavyzdžiui, „l“, „s“ ar „r“.} \\
\texttt{spėjimas: aloud} \\
            }
        }
    }
    & & \\ \\

    \theutterance \stepcounter{utterance}  
    & & & \multicolumn{2}{p{0.3\linewidth}}{
        \cellcolor[rgb]{0.9,0.9,0.9}{
            \makecell[{{p{\linewidth}}}]{
                \texttt{\tiny{[GM$|$GM]}}
                \texttt{Error: The guessed word is not a valid word for this game.} \\
            }
        }
    }
    & & \\ \\

    \theutterance \stepcounter{utterance}  
    & & & \multicolumn{4}{p{0.6\linewidth}}{
        \cellcolor[rgb]{0.9,0.9,0.9}{
            \makecell[{{p{\linewidth}}}]{
                \texttt{\tiny{[P1$\langle$GM]}}
                \texttt{Jūsų spėjimas nėra tinkamas žodis šiam žaidimui. Pabandykite dar kartą.} \\
\\ 
\texttt{Atsakymą pateikite tik šiuo formatu.} \\
\texttt{paaiškinimas: detalės} \\
\texttt{spėjimas: žodis} \\
            }
        }
    }
     \\ \\

    \theutterance \stepcounter{utterance}  
    & \multicolumn{4}{p{0.6\linewidth}}{
        \cellcolor[rgb]{0.561,0.737,0.561}{
            \makecell[{{p{\linewidth}}}]{
                \texttt{\tiny{[P1$\rangle$GM]}}
                \texttt{paaiškinimas: Kadangi „aloud“ buvo priimtas kaip netinkamas žodis, grįžtu prie ankstesnių duomenų: žodis prasideda su „a“ ir baigiasi su „o“. Raidė „u“ yra žodžio dalyje, bet ne trečioje pozicijoje. „D“ ir „i“ yra neteisingos. Išbandysiu kitą galimą kombinaciją su „u“ raidė ne trečioje vietoje, naudojant kitas dažnas raides, pavyzdžiui, „l“ arba „s“.} \\
\texttt{spėjimas: autos} \\
            }
        }
    }
    & & \\ \\

    \theutterance \stepcounter{utterance}  
    & & & \multicolumn{2}{p{0.3\linewidth}}{
        \cellcolor[rgb]{0.9,0.9,0.9}{
            \makecell[{{p{\linewidth}}}]{
                \texttt{\tiny{[GM$|$GM]}}
                \texttt{attempts = 2} \\
\texttt{target\_word = akmuo} \\
\texttt{guess = autos} \\
\texttt{guess\_feedback = a<green> u<yellow> t<red> o<yellow> s<red>} \\
            }
        }
    }
    & & \\ \\

    \theutterance \stepcounter{utterance}  
    & & & \multicolumn{4}{p{0.6\linewidth}}{
        \cellcolor[rgb]{0.9,0.9,0.9}{
            \makecell[{{p{\linewidth}}}]{
                \texttt{\tiny{[P1$\langle$GM]}}
                \texttt{spėjimo\_atsiliepimas: a<žalia> u<geltona> t<raudona> o<geltona> s<raudona>} \\
\\ 
\texttt{Atsakymą pateikite tik šiuo formatu.} \\
\texttt{paaiškinimas: detalės} \\
\texttt{spėjimas: žodis} \\
            }
        }
    }
     \\ \\

    \theutterance \stepcounter{utterance}  
    & \multicolumn{4}{p{0.6\linewidth}}{
        \cellcolor[rgb]{0.561,0.737,0.561}{
            \makecell[{{p{\linewidth}}}]{
                \texttt{\tiny{[P1$\rangle$GM]}}
                \texttt{paaiškinimas: Žinau, kad žodis prasideda su „a“. Raidė „u“ yra žodžio dalyje, bet ne antroje ir ne trečioje pozicijoje. Raidė „o“ yra žodžio dalyje, bet ne ketvirtoje pozicijoje. Raidės „t“ ir „s“ yra neteisingos. Bandysiu suformuoti žodį, kuriame „u“ būtų ketvirtoje pozicijoje, o „o“ – trečioje, naudojant naujas raides, pavyzdžiui, „m“, „n“ ar „l“.} \\
\texttt{spėjimas: amuso} \\
            }
        }
    }
    & & \\ \\

    \theutterance \stepcounter{utterance}  
    & & & \multicolumn{2}{p{0.3\linewidth}}{
        \cellcolor[rgb]{0.9,0.9,0.9}{
            \makecell[{{p{\linewidth}}}]{
                \texttt{\tiny{[GM$|$GM]}}
                \texttt{Error: The guessed word is not a valid word for this game.} \\
            }
        }
    }
    & & \\ \\

    \theutterance \stepcounter{utterance}  
    & & & \multicolumn{4}{p{0.6\linewidth}}{
        \cellcolor[rgb]{0.9,0.9,0.9}{
            \makecell[{{p{\linewidth}}}]{
                \texttt{\tiny{[P1$\langle$GM]}}
                \texttt{Jūsų spėjimas nėra tinkamas žodis šiam žaidimui. Pabandykite dar kartą.} \\
\\ 
\texttt{Atsakymą pateikite tik šiuo formatu.} \\
\texttt{paaiškinimas: detalės} \\
\texttt{spėjimas: žodis} \\
            }
        }
    }
     \\ \\

    \theutterance \stepcounter{utterance}  
    & \multicolumn{4}{p{0.6\linewidth}}{
        \cellcolor[rgb]{0.561,0.737,0.561}{
            \makecell[{{p{\linewidth}}}]{
                \texttt{\tiny{[P1$\rangle$GM]}}
                \texttt{paaiškinimas: Kadangi „amuso“ buvo atmetamas, ieškau realaus anglų kalbos žodžio. Žinau, kad žodis prasideda su „a“. Raidė „u“ yra žodžio dalyje, bet ne antrosios ar trečiosios pozicijos. Raidė „o“ yra žodžio dalyje, bet ne ketvirtosios pozicijos. Raidės „t“, „s“, „d“ ir „i“ yra neteisingos. Bandysiu žodį, kuriame „u“ būtų penktoje pozicijoje, o „o“ – antroje, naudojant naujas raides, pavyzdžiui, „l“ ir „n“.} \\
\texttt{spėjimas: aloud} \\
            }
        }
    }
    & & \\ \\

    \theutterance \stepcounter{utterance}  
    & & & \multicolumn{2}{p{0.3\linewidth}}{
        \cellcolor[rgb]{0.9,0.9,0.9}{
            \makecell[{{p{\linewidth}}}]{
                \texttt{\tiny{[GM$|$GM]}}
                \texttt{Error: The guessed word is not a valid word for this game.} \\
            }
        }
    }
    & & \\ \\

    \theutterance \stepcounter{utterance}  
    & & & \multicolumn{4}{p{0.6\linewidth}}{
        \cellcolor[rgb]{0.9,0.9,0.9}{
            \makecell[{{p{\linewidth}}}]{
                \texttt{\tiny{[P1$\langle$GM]}}
                \texttt{Jūsų spėjimas nėra tinkamas žodis šiam žaidimui. Pabandykite dar kartą.} \\
\\ 
\texttt{Atsakymą pateikite tik šiuo formatu.} \\
\texttt{paaiškinimas: detalės} \\
\texttt{spėjimas: žodis} \\
            }
        }
    }
     \\ \\

    \theutterance \stepcounter{utterance}  
    & \multicolumn{4}{p{0.6\linewidth}}{
        \cellcolor[rgb]{0.561,0.737,0.561}{
            \makecell[{{p{\linewidth}}}]{
                \texttt{\tiny{[P1$\rangle$GM]}}
                \texttt{paaiškinimas: Atsitiko klaida, nes „aloud“ jau bandiau ir jis buvo atmetamas. Reikia suformuoti realų žodį. Žinau, kad žodis prasideda su „a“. Raidė „u“ yra žodžio dalyje, bet ne 2{-}oje, 3{-}ioje ar 4{-}oje pozicijoje (iš „autos“ ir „audio“). Raidė „o“ yra žodžio dalyje, bet ne 4{-}oje pozicijoje. Raidės „d“, „i“, „t“, „s“ yra neteisingos. Jei „u“ nėra 2, 3 ar 4 vietoje, ji turi būti 5 vietoje. Jei „o“ nėra 4 vietoje, bandysiu ją 2 vietoje. Išbandysiu žodį su „a“, „o“ ir „u“.} \\
\texttt{spėjimas: adieu} \\
            }
        }
    }
    & & \\ \\

    \theutterance \stepcounter{utterance}  
    & & & \multicolumn{2}{p{0.3\linewidth}}{
        \cellcolor[rgb]{0.9,0.9,0.9}{
            \makecell[{{p{\linewidth}}}]{
                \texttt{\tiny{[GM$|$GM]}}
                \texttt{Error: The guessed word is not a valid word for this game.} \\
            }
        }
    }
    & & \\ \\

    \theutterance \stepcounter{utterance}  
    & & & \multicolumn{4}{p{0.6\linewidth}}{
        \cellcolor[rgb]{0.9,0.9,0.9}{
            \makecell[{{p{\linewidth}}}]{
                \texttt{\tiny{[P1$\langle$GM]}}
                \texttt{Jūsų spėjimas nėra tinkamas žodis šiam žaidimui. Pabandykite dar kartą.} \\
\\ 
\texttt{Atsakymą pateikite tik šiuo formatu.} \\
\texttt{paaiškinimas: detalės} \\
\texttt{spėjimas: žodis} \\
            }
        }
    }
     \\ \\

    \theutterance \stepcounter{utterance}  
    & \multicolumn{4}{p{0.6\linewidth}}{
        \cellcolor[rgb]{0.561,0.737,0.561}{
            \makecell[{{p{\linewidth}}}]{
                \texttt{\tiny{[P1$\rangle$GM]}}
                \texttt{paaiškinimas: Atmetami žodžiai „aloud“ ir „adieu“ rodo, kad turi būti atsargiai su raidėmis „d“ ir „i“. Žinau, kad žodis prasideda su „a“. Raidė „u“ yra žodžio dalyje, bet ne 2, 3 ar 4 pozicijoje. Raidė „o“ yra žodžio dalyje, bet ne 4 pozicijoje. Jei „u“ yra 5 pozicijoje, o „o“ yra 2 pozicijoje, bandysiu žodį su naujomis raidėmis „m“, „n“ ar „l“.} \\
\texttt{spėjimas: amoun} \\
            }
        }
    }
    & & \\ \\

    \theutterance \stepcounter{utterance}  
    & & & \multicolumn{2}{p{0.3\linewidth}}{
        \cellcolor[rgb]{0.9,0.9,0.9}{
            \makecell[{{p{\linewidth}}}]{
                \texttt{\tiny{[GM$|$GM]}}
                \texttt{Error: The guessed word is not a valid word for this game.} \\
            }
        }
    }
    & & \\ \\

    \theutterance \stepcounter{utterance}  
    & & & \multicolumn{4}{p{0.6\linewidth}}{
        \cellcolor[rgb]{0.9,0.9,0.9}{
            \makecell[{{p{\linewidth}}}]{
                \texttt{\tiny{[P1$\langle$GM]}}
                \texttt{Jūsų spėjimas nėra tinkamas žodis šiam žaidimui. Pabandykite dar kartą.} \\
\\ 
\texttt{Atsakymą pateikite tik šiuo formatu.} \\
\texttt{paaiškinimas: detalės} \\
\texttt{spėjimas: žodis} \\
            }
        }
    }
     \\ \\

    \theutterance \stepcounter{utterance}  
    & \multicolumn{4}{p{0.6\linewidth}}{
        \cellcolor[rgb]{0.561,0.737,0.561}{
            \makecell[{{p{\linewidth}}}]{
                \texttt{\tiny{[P1$\rangle$GM]}}
                \texttt{paaiškinimas: Kadangi ankstesni spėjimai buvo atmetami kaip netinkami žodžiai, turiu rasti realų anglų kalbos žodį. Žinau, kad žodis prasideda su „a“. Raidė „u“ yra žodžio dalyje, bet ne 2, 3 ar 4 pozicijoje. Raidė „o“ yra žodžio dalyje, bet ne 4 pozicijoje. Jei „u“ yra 5 pozicijoje, o „o“ yra 2 pozicijoje, bandysiu žodį „anode“, tačiau „d“ jau buvo nurodyta kaip neteisinga. Bandysiu žodį „abono“, tačiau tai nėra anglų žodis. Išbandysiu žodį „anode“, bet „d“ yra neteisinga. Bandysiu „aloud“ vėl, bet tai neįvyks. Išbandysiu „anode“ variantą be „d“: „anole“.} \\
\texttt{spėjimas: anole} \\
            }
        }
    }
    & & \\ \\

    \theutterance \stepcounter{utterance}  
    & & & \multicolumn{2}{p{0.3\linewidth}}{
        \cellcolor[rgb]{0.9,0.9,0.9}{
            \makecell[{{p{\linewidth}}}]{
                \texttt{\tiny{[GM$|$GM]}}
                \texttt{Error: The guessed word is not a valid word for this game.} \\
            }
        }
    }
    & & \\ \\

    \theutterance \stepcounter{utterance}  
    & & & \multicolumn{4}{p{0.6\linewidth}}{
        \cellcolor[rgb]{0.9,0.9,0.9}{
            \makecell[{{p{\linewidth}}}]{
                \texttt{\tiny{[P1$\langle$GM]}}
                \texttt{Jūsų spėjimas nėra tinkamas žodis šiam žaidimui. Pabandykite dar kartą.} \\
\\ 
\texttt{Atsakymą pateikite tik šiuo formatu.} \\
\texttt{paaiškinimas: detalės} \\
\texttt{spėjimas: žodis} \\
            }
        }
    }
     \\ \\

    \theutterance \stepcounter{utterance}  
    & \multicolumn{4}{p{0.6\linewidth}}{
        \cellcolor[rgb]{0.561,0.737,0.561}{
            \makecell[{{p{\linewidth}}}]{
                \texttt{\tiny{[P1$\rangle$GM]}}
                \texttt{paaiškinimas: Nuolat gaunamas pranešimas, kad žodžiai nėra tinkami, rodo, kad mano bandymai suformuoti realius anglų kalbos žodžius su nurodyta struktūra yra klaidingi. Reikia grįžti prie bazinių faktų: žodis prasideda su „a“, raidė „u“ yra žodžio dalyje (bet ne 2, 3 ar 4 pozicijoje), raidė „o“ yra žodžio dalyje (bet ne 4 pozicijoje). Jei „u“ yra 5 pozicijoje, o „o“ yra 2 pozicijoje, bandysiu žodį „amour“, kuris yra realus anglų kalbos žodis (išsamus meilės žodis).} \\
\texttt{spėjimas: amour} \\
            }
        }
    }
    & & \\ \\

    \theutterance \stepcounter{utterance}  
    & & & \multicolumn{2}{p{0.3\linewidth}}{
        \cellcolor[rgb]{0.9,0.9,0.9}{
            \makecell[{{p{\linewidth}}}]{
                \texttt{\tiny{[GM$|$GM]}}
                \texttt{Error: The guessed word is not a valid word for this game.} \\
            }
        }
    }
    & & \\ \\

    \theutterance \stepcounter{utterance}  
    & & & \multicolumn{2}{p{0.3\linewidth}}{
        \cellcolor[rgb]{0.9,0.9,0.9}{
            \makecell[{{p{\linewidth}}}]{
                \texttt{\tiny{[GM$|$GM]}}
                \texttt{game\_result = ABORT} \\
            }
        }
    }
    & & \\ \\

\end{supertabular}
}

\end{document}
